\documentclass[12pt]{article}
\usepackage[T1]{fontenc}
\usepackage[utf8]{inputenc}
\usepackage{lmodern}
\usepackage{textcomp}
\usepackage{enumitem}
\usepackage{geometry}
\geometry{margin=1in}
\begin{document}
\begin{center}
Answer Key - Philosophy - Graduate Level Assessment 3 \
\small (Answers are ordered exactly as in the question paper.)
\end{center}

\section*{Multiple Choice}
\begin{enumerate}[label=\arabic*.]
\item \textbf{Answer:} A
\\ \textit{Options:} A) failure to abide by a contract; B) disregard for societal norms; C) acting against the welfare of others; D) disobedience to parental authority; E) disobedience to God's law
\\ \textit{Note:} According to Hobbes, injustice arises when individuals fail to fulfill their contractual obligations, as he believes that the social contract is fundamental to maintaining order and justice in society.
\item \textbf{Answer:} A
\\ \textit{Options:} A) one's own will.; B) one's own emotions.; C) God.; D) Nature.; E) the universe.
\\ \textit{Note:} Kant argues that the moral law is an expression of rational will, meaning that individuals have the capacity to legislate moral laws for themselves through their own rational deliberation, making one's own will the source of moral authority.
\item \textbf{Answer:} A
\\ \textit{Options:} A) a person who donates to charity regularly.; B) all of the above.; C) a person who lived in Nazi Germany in 1940.; D) a person who drives drunk and hits a child.; E) none of the above.
\\ \textit{Note:} The correct answer highlights that a person who donates to charity regularly is subject to moral luck because their moral character is judged based on their actions, which can be influenced by external circumstances beyond their control, illustrating the complexities of evaluating moral responsibility.
\item \textbf{Answer:} A
\\ \textit{Options:} A) loss of human diversity; B) the effects of engineering children to have positional goods; C) loss of equality of opportunity; D) increased disparity between the rich and the poor
\\ \textit{Note:} Singer believes that while the loss of human diversity is concerning, it is less serious than other potential social problems associated with a genetic supermarket, such as the commodification of human life and the ethical implications of genetic selection, which pose more immediate threats to human dignity and equality.
\item \textbf{Answer:} B
\\ \textit{Options:} A) quantum physics; B) metaphysics; C) ontology; D) phenomenology; E) cosmology
\\ \textit{Note:} Metaphysics is the branch of philosophy that explores the fundamental nature of reality, including the relationship between mind and matter, substance and attribute, and potentiality and actuality.
\end{enumerate}

\section*{True / False}
\begin{enumerate}[label=\arabic*.]
\setcounter{enumi}{5}
\item \textbf{Answer:} False
\\ \textit{Options:} A) True; B) False
\\ \textit{Note:} Just war theory asserts that moral considerations are central to the justification and conduct of war, emphasizing the importance of ethical principles in determining when it is permissible to engage in conflict and how to engage in it justly.
\item \textbf{Answer:} True
\\ \textit{Options:} A) True; B) False
\\ \textit{Note:} Greek philosophers began to reject the anthropomorphic view of the divine in the late Sixth Century BCE as they sought more abstract and rational explanations for the nature of existence and the cosmos, moving towards a conception of divinity that emphasized an impersonal, metaphysical reality rather than human-like deities.
\item \textbf{Answer:} False
\\ \textit{Options:} A) True; B) False
\\ \textit{Note:} Denying the antecedent in a conditional syllogism does not lead to affirming the consequent because it violates the logical structure of conditional statements, where the truth of the consequent cannot be confirmed solely based on the denial of the antecedent.
\item \textbf{Answer:} False
\\ \textit{Options:} A) True; B) False
\\ \textit{Note:} The Tripitaka, also known as the Pali Canon, actually refers to the "Three Baskets" of teachings that encompass the Vinaya Pitaka, Sutta Pitaka, and Abhidhamma Pitaka, rather than the 'Three teachings' of Buddhist philosophy.
\item \textbf{Answer:} False
\\ \textit{Options:} A) True; B) False
\\ \textit{Note:} The Oral Torah, in its written form, is referred to as the Mishnah, while the Talmud serves as a comprehensive compilation of rabbinic discussions and interpretations that expand upon the Mishnah.
\end{enumerate}

\section*{Long Form Response}
\begin{enumerate}[label=\arabic*.]
\setcounter{enumi}{10}
\item \textbf{Answer:} Aldo Leopold argues that ecological issues cannot be reduced to mere economic considerations, as they encompass a broader moral and ethical dimension. He emphasizes that our relationship with the natural environment involves responsibilities that transcend profit and utility, calling for a land ethic that recognizes the intrinsic value of ecosystems. Leopold asserts that environmental degradation results not only from economic exploitation but from a fundamental disconnect between humanity and the natural world. By framing ecological concerns as ethical imperatives, he advocates for a holistic understanding of our role within the biotic community, urging a shift towards stewardship rather than ownership. Thus, the challenges we face require a reevaluation of our values, prioritizing ecological integrity over short-term economic gains.
\\ \textit{Note:} correct because it captures Leopold's belief that ecological issues involve ethical responsibilities and intrinsic values that go beyond economic interests, highlighting the need for a land ethic to address the disconnect between humanity and the environment.
\item \textbf{Answer:} Stevenson confronts the objection concerning the lack of objective, a priori goodness in his moral theory by expressing a fundamental misunderstanding of such a property. He argues that the concept of objective goodness is elusive and does not align with his emphasis on the emotive and expressive dimensions of moral language. By asserting that he does not comprehend the notion of an a priori goodness, Stevenson highlights the nature of moral properties as inherently subjective and contingent upon human feelings and social contexts. This response underscores his view that moral evaluations are not grounded in fixed, objective standards but are instead shaped by the interplay of human sentiments and social interactions. Thus, Stevenson's position invites a re-examination of how moral properties are conceived, steering the discourse away from traditional notions of objectivity in ethics.
\\ \textit{Note:} The answer correctly identifies Stevenson's challenge to the notion of objective, a priori goodness by emphasizing his focus on the emotive aspects of moral language, thereby asserting that moral properties are fundamentally subjective and context-dependent.
\end{enumerate}

\vfill
\noindent\textit{End of Answer Key}
\end{document}